% Part: first-order-logic
% Chapter: natural-deduction
% Section: quantifier-rules

\documentclass[../../../include/open-logic-section]{subfiles}

\begin{document}

\olfileid{fol}{ntd}{qrl}

\olsection{संख्यापकांचे नियम}

\subsection{$\lforall$ साठीचे नियम}

\begin{defish}
\AxiomC{$!A(a)$}
\RightLabel{\Intro{\lforall}}
\UnaryInfC{$\lforall[x][\Atom{!A}{x}]$}
\DisplayProof
\hfill
\AxiomC{$\lforall[x][\Atom{!A}{x}]$}
\RightLabel{\Elim{\lforall}}
\UnaryInfC{$!A(t)$}
\DisplayProof
\end{defish}

$\lforall$ च्या नियमांमध्ये, $t$ हे बंद पद आहे (म्हणजे, कोणतेही चर
नसलेले पद), आणि $a$~हा असा !!a{constant} आहे की तो
निष्कर्ष~$\lforall[x][!A(x)]$ मध्ये किंवा आधारविधान~$!A(a)$ ने समाप्त
होणाऱ्या !!{derivation} मध्ये !!{undischarged} असलेल्या कोणत्याही
गृहीतकात येत नाही. आपण $a$ ला \Intro{\lforall} अनुमानाचा
\emph{आयगेन चर} म्हणतो.\footnote{वरील नियमात $a$ हा अचर असला तरी आपण
``आयगेन चर'' ही संज्ञा वापरतो. यामागे ऐतिहासिक कारणे आहेत.}

\subsection{$\lexists$ साठीचे नियम}

\begin{defish}
\AxiomC{$\Atom{!A}{t}$}
\RightLabel{\Intro{\lexists}}
\UnaryInfC{$\lexists[x][\Atom{!A}{x}]$}
\DisplayProof
\hfill
\AxiomC{$\lexists[x][\Atom{!A}{x}]$}
\AxiomC{[$\Atom{!A}{a}$]$^n$}
\DeduceC{$!C$}
\DischargeRule{\Elim{\lexists}}{n}
\BinaryInfC{$!C$}
\DisplayProof
\end{defish}

पुन्हा, $t$ हे बंद पद आहे आणि $a$ हा असा !!a{constant} आहे की तो
आधारविधान $\lexists[x][!A(x)]$ मध्ये, निष्कर्ष~$!C$ मध्ये किंवा दोन
आधारविधानांनी समाप्त होणाऱ्या !!{derivation}s मधील कोणत्याही
!!{undischarged} गृहीतकात ($!A(a)$ ही गृहीतके वगळता) येत नाही. आपण
$a$ ला \Elim{\lexists} अनुमानाचा \emph{आयगेन चर} म्हणतो.

\Intro{\lforall} किंवा \Elim{\lexists} अनुमानांच्या आधारविधानांकडे
नेणाऱ्या !!{derivation}s मध्ये आयगेन चर आधारविधानांत किंवा कोणत्याही
!!{undischarged} गृहीतकात येत नसावा, या अटीला \emph{आयगेन चराची अट}
म्हणतात.

\begin{explain}
ही प्रथा आठवा: $!A$ हे असे !!a{formula} असेल की त्यात
!!{variable}~$x$ मुक्त असेल, तर हे आपण~$!A(x)$ असे लिहून दर्शवतो.
त्याच संदर्भात, मग $!A(t)$ हा~$\Subst{!A}{t}{x}$ चा संक्षेप असतो.
त्यामुळे $\Intro\lexists$ नियम पुढीलप्रमाणेही लिहिता येईल:
\begin{prooftree}
\AxiomC{$\Subst{!A}{t}{x}$}
\RightLabel{\Intro{\lexists}}
\UnaryInfC{$\lexists[x][!A]$}
\end{prooftree}
लक्षात घ्या की~$!A$ मध्ये $t$ आधीपासून येऊ शकते; उदा., $!A$ हे
$\Atom{\Obj P}{t,x}$ असू शकते. म्हणून,~$\Atom{\Obj P}{t,t}$ पासून
$\lexists[x][\Atom{\Obj P}{t,x}]$ निष्पन्न करणे हे~$\Intro\lexists$
चे योग्य उपयोजन आहे---आधारविधानात~$t$ जेथे येते त्यांपैकी एक किंवा
अधिक ठिकाणी त्याऐवजी बद्ध !!{variable}~$x$ घालता येतो; सर्वच ठिकाणी
असा ``बदल'' करणे आवश्यक नाही. परंतु $\Intro\lforall$ आणि
$\Elim\lexists$ मधील आयगेन चराच्या अटीनुसार !!{constant}~$a$ हा~$!A$
मध्ये येता कामा नये. म्हणून, $\Intro\lforall$ वापरून
$\Atom{\Obj P}{a,a}$ पासून $\lforall[x][\Atom{\Obj P}{a,x}]$ योग्य
रीतीने निष्पन्न करता येत नाही.
\end{explain}

\begin{explain}
\Intro{\lexists} आणि \Elim{\lforall} मध्ये पद~$t$ बंद
असण्याव्यतिरिक्त त्यावर आणखी कोणतेही निर्बंध नाहीत; म्हणून इतर कोणत्याही
अटीची काळजी करावी लागत नाही. याउलट, \Elim{\lexists} आणि
\Intro{\lforall} नियमांमध्ये आयगेन चराच्या अटीनुसार !!{constant}~$a$ हा
निष्कर्षात किंवा कोणत्याही !!{undischarged} गृहीतकात कोठेही येता कामा नये.
ही पद्धत निर्दोष आहे, म्हणजे ती केवळ अशा !!{undischarged} गृहीतकांपासूनच
!!{sentence}s !!{derive}s करते की ज्यांच्यापासून ती तार्किकरीत्या निष्पन्न
होतात, याची खात्री करण्यासाठी ही अट आवश्यक आहे. ही अट नसती, तर पुढील
निष्पत्तीला परवानगी मिळाली असती:
\begin{prooftree}
  \AxiomC{$\lexists[x][!A(x)]$}
  \AxiomC{$\Discharge{!A(a)}{1}$}
  \RightLabel{*\Intro{\lforall}}
  \UnaryInfC{$\lforall[x][!A(x)]$}
  \RightLabel{\Elim{\lexists}}
  \BinaryInfC{$\lforall[x][!A(x)]$}
\end{prooftree}
परंतु, $\lexists[x][!A(x)] \Entails/ \lforall[x][!A(x)]$.

संख्यापकांच्या विलोपन नियमांमध्ये !!{variable}s च्या जागी केवळ बंद पदे
ठेवण्याची परवानगी असल्यामुळे, !!{sentence}s च्या संचापासून निष्पन्न करता
येणारे कोणतेही !!{formula} हे स्वतः !!a{sentence} असते, हे निष्पन्न होते.
\end{explain}


\end{document}
